\section*{Bab \ref{ch:g12:diffeq} --- Persamaan Diferensial}

\begin{solution}{exo:g12:diffeq:1}
\emph{(a)} $y = C\eu^{3x}$; $y(0) = 2$ memberi $y = 2\eu^{3x}$.

\emph{(b)} $y' = -\frac12 y$, jadi $y = C\eu^{-x/2}$; $y(0) = -1$ memberi
$y = -\eu^{-x/2}$.

\emph{(c)} Kesetimbangannya $y_p = 5$; $y = C\eu^{-x} + 5$; $y(0) = 0$ memberi
$C = -5$, jadi $y = 5\left(1 - \eu^{-x}\right)$.
\end{solution}

\begin{solution}{exo:g12:diffeq:2}
$N' = kN$, jadi $N(t) = N_0 \eu^{kt}$. Berlipat dua dalam $3$ jam berarti
$\eu^{3k} = 2$, \emph{yaitu}
\[
k = \frac{\ln 2}{3} \approx 0.231\ \text{jam}^{-1}.
\]
\end{solution}

\begin{solution}{exo:g12:diffeq:3}
$y_p'(x) = \eu^x + x\eu^x = y_p(x) + \eu^x$: jadi $y_p$ satu penyelesaian khusus
(inilah kasus beresonansi pada \cref{met:g12:diffeq:particular}). Menurut
\cref{prop:g12:diffeq:general}, penyelesaiannya adalah
$y = C\eu^{x} + x\eu^{x} = (C + x)\,\eu^x$, $C \in \R$.
\end{solution}

\begin{solution}{exo:g12:diffeq:4}
Cobalah $y_p = \beta\eu^x$: $\beta\eu^x = -2\beta\eu^x + \eu^x$ memberi
$3\beta = 1$, jadi $y_p = \frac13\eu^x$. Penyelesaian umumnya
$y = C\eu^{-2x} + \frac13\eu^{x}$; syarat $y(0) = 1$ memberi
$C = \frac23$:
\[
y(x) = \frac{2}{3}\,\eu^{-2x} + \frac{1}{3}\,\eu^{x}.
\]
\end{solution}

\begin{solution}{exo:g12:diffeq:5}
$N(t) = N_0\,\eu^{-\lambda t}$ dengan $\lambda = \frac{\ln 2}{5730}$
(\cref{ex:g12:exp:halflife}). Kita selesaikan $\eu^{-\lambda t} = 0.6$:
\[
t = \frac{\ln(1/0.6)}{\lambda} = 5730\,\frac{\ln(5/3)}{\ln 2}
\approx 5730 \times \frac{0.5108}{0.6931} \approx 4220 \text{ tahun}.
\]
\end{solution}

\begin{solution}{exo:g12:diffeq:6}
\emph{1.} Garamnya mengalir masuk dengan laju $0.2 \times 5 = 1$ kg/menit. Aliran
keluarnya membawa kepekatan $\frac{m}{100}$ kg/L pada $5$ L/menit, \emph{yaitu}
$\frac{m}{20}$ kg/menit. Jadi $m' = 1 - \frac{m}{20}$.

\emph{2.} Kesetimbangannya $m_p = 20$; $m(t) = 20 + C\eu^{-t/20}$, dan
$m(0) = 0$ memberi $C = -20$:
\[
m(t) = 20\left(1 - \eu^{-t/20}\right) \xrightarrow[t\to+\infty]{} 20 \text{ kg}.
\]
Dalam jangka panjang kepekatan tangkinya sama dengan kepekatan air garam yang
masuk: $0.2$ kg/L $\times$ $100$ L $= 20$ kg.
\end{solution}

\begin{solution}{exo:g12:diffeq:7}
\emph{1.} $v' = g - \frac{k}{m} v$ dengan $\frac km = \frac{16}{80} = 0.2$.
Kesetimbangannya $v_\infty = \frac{mg}{k} = \frac{9.8}{0.2} = 49$ m/s;
$v(t) = 49 + C\eu^{-0.2t}$, dan $v(0) = 0$ memberi
\[
v(t) = 49\left(1 - \eu^{-0.2 t}\right).
\]

\emph{2.} Kecepatan terminalnya $49$ m/s ($\approx 176$ km/jam). Kita
inginkan $1 - \eu^{-0.2t} = 0.95$, \emph{yaitu} $\eu^{-0.2t} = 0.05$:
\[
t = \frac{\ln 20}{0.2} \approx \frac{3.00}{0.2} \approx 15 \text{ sekon}.
\]
\end{solution}

\begin{solution}{exo:g12:diffeq:8}
\emph{1.} $z = \frac1y$ memberi $z' = -\frac{y'}{y^2}
= -\frac{y(1-y)}{y^2} = -\frac{1-y}{y} = -\frac1y + 1 = -z + 1$.

\emph{2.} Kesetimbangannya $z_p = 1$, jadi $z(t) = 1 + C\eu^{-t}$. Dari
$y(0) = \frac{1}{10}$ diperoleh $z(0) = 10$, jadi $C = 9$ dan
\[
y(t) = \frac{1}{1 + 9\,\eu^{-t}} .
\]

\emph{3.} Ketika $t \to +\infty$, $9\eu^{-t} \to 0$ dan $y(t) \to 1$. Kurvanya
adalah kurva logistik berbentuk S yang klasik (\emph{sigmoid}): pertumbuhannya
lambat mula-mula ($y$ kecil, $y' \approx y$), tercepat ketika $y = \frac12$
(ketika $y' = y(1-y)$ maksimal), lalu menjenuh menuju daya dukungnya
yang bernilai $1$.
\end{solution}

\begin{solution}{exo:g12:diffeq:9}
Menurut \cref{thm:g12:diffeq:affine}, $y(x) = C\eu^{ax} - \frac ba$. Jadi
\[
u_{n+1} = C\eu^{a(n+1)} - \frac ba
= \eu^{a}\left(C\eu^{an} - \frac ba\right) + \frac ba\left(\eu^a - 1\right)
= q\,u_n + r
\]
dengan $q = \eu^{a}$ dan $r = \frac{b}{a}\left(\eu^{a} - 1\right)$.
Syarat $\abs q < 1$ berarti $\eu^a < 1$, \emph{yaitu} $a < 0$: tepat
syarat yang membuat penyelesaian \omterm{def:g12:diffeq:ode}{persamaan diferensialnya}
menuju kesetimbangan $-\frac ba$ --- dan memang \omterm{pb:g10:functions:1}{titik tetap}
rekurensinya adalah $\frac{r}{1 - q} = -\frac{b}{a}$.
\end{solution}

\begin{solution}{pb:g12:diffeq:1}
\textbf{1.} $y = 2\eu^{3t}$. Untuk yang kedua: kesetimbangannya $3$,
jadi $y = 3 + C\eu^{-2t}$ dengan $y(0) = 0$, sehingga $C = -3$ dan
$y = 3\left(1 - \eu^{-2t}\right)$.

\textbf{2.} $\left(y\eu^{-at}\right)' = y'\eu^{-at} -
ay\eu^{-at} = (y' - ay)\eu^{-at} = 0$: jadi hasil kalinya sebuah
tetapan $C$, sehingga $y = C\eu^{at}$ --- tak ada penyelesaian yang lolos.

\textbf{3.} Kesetimbangannya: $y \equiv 3$. Setiap penyelesaian lainnya berupa
$3 + C\eu^{-2t}$ dengan $C \neq 0$: \omterm{thm:g12:exp:existence}{fungsi eksponennya} mati dan
penyelesaiannya meluncur ke $3$ --- sebuah \omterm{pb:g10:functions:1}{titik tetap} yang mantap, kembaran
sinambung dari \omterm{pb:g10:functions:1}{titik tetap} rekurensi pada
\cref{pb:g11:seq:1}.

\textbf{4.} $y = y_0\eu^{-t/\tau}$ menjadi separuhnya ketika
$\eu^{-t/\tau} = \frac12$, jadi $t = \tau\ln 2$.

\textbf{5.} Semua penyelesaiannya adalah geseran tegak dari peluruhan
menuju $3$: yang bermula di atasnya jatuh ke $3$, yang di bawahnya
naik ke $3$, dan penyelesaian tetap $3$ diam saja --- sebuah
corong di sekeliling kesetimbangannya.

\textbf{6.} $\lambda = \frac{\ln 2}{5730} \approx 1.21 \times
10^{-4}$ tiap tahunnya.

\textbf{7.} $\eu^{-\lambda t} = 0.2$, jadi
$t = \frac{\ln 5}{\lambda} \approx 13\,300$ tahun.

\textbf{8.} $t = \frac{\ln(1/0.15)}{\lambda} \approx 15\,700$
tahun: banteng Lascaux itu berasal dari akhir zaman es.

\textbf{9.} $t = \frac{\ln(1/0.53)}{\lambda} \approx 5\,250$
tahun: hanya selisih satu umur manusia dari angka para arkeolog ---
jamnya bekerja.

\textbf{10.} Sepuluh waktu paruh menyisakan $2^{-10} \approx 0.1\,\%$,
sudah di tepi keterukuran. Sesudah $65$ juta tahun
pecahannya menjadi $2^{-65\,000\,000/5\,730} \approx 2^{-11\,344}$:
tak satu atom pun dari persediaan asalnya tersisa di fosil mana pun ---
dinosaurus ditanggali dengan jam yang lebih lambat (kalium--argon dan
kerabatnya).

\textbf{11.} Nilai $52\,\%$ memberi $\approx 5\,405$ tahun, dan
$54\,\%$ memberi $\approx 5\,095$: jadi $\pm 1\,\%$ milik kimia
menjadi kira-kira $\pm 155$ tahun milik sejarah --- ketelitian di
laboratorium adalah ketelitian pada label museumnya.

\textbf{12.} $z = T - T_a$ memenuhi $z' = -kz$, jadi
$z = (T_0 - T_a)\eu^{-kt}$, dan $T = T_a + z$.

\textbf{13.} Jurang awalnya $T_0 - T_a = 70$; sesudah lima
menit jurangnya $70 - 20 = 50$, jadi $50 = 70\eu^{-5k}$ dan
$k = \frac{\ln(70/50)}{5} \approx 0.067$ tiap menitnya. Enak diminum:
$55 - 20 = 35 = 70\eu^{-kt}$, jadi
$t = \frac{\ln 2}{k} \approx 10.3$ menit.

\textbf{14.} Laju kehilangan panasnya sebanding dengan jurang
$T - T_a$: kopi yang mendidih membuang panas paling cepat. Menambahkan
susunya \emph{sekarang} langsung memangkas jurangnya, sehingga panas yang
hilang selama sepuluh menit itu lebih sedikit; menambahkannya di akhir membiarkan
kopinya mendingin dengan laju penuh lebih dahulu. Untuk cangkir terhangat: susunya
dahulu. (Tuan rumah yang tidak sabar membalik urutan termodinamikanya.)

\textbf{15.} Jurang terhadap ruangannya: pada tengah malam $10$, satu jam
kemudian $8$, jadi $\eu^{-k} = 0.8$ dan
$k = \ln\frac{10}{8} \approx 0.223$ tiap jamnya. Pada saat mati jurangnya
adalah $17$; jurang itu meluruh menjadi $10$ pada tengah malam, sehingga
$s = \frac{\ln(17/10)}{k} \approx 2.4$ jam. Saat kematiannya:
kira-kira pukul $21{:}40$.

\textbf{16.} Ruangan yang dipanaskan atau berangin ($T_a$ tidak tetap)
membengkokkan jamnya ke dua arah; pakaian atau massa tubuh mengubah $k$
(tabel dokter forensik menyesuaikannya) --- $k$ yang keliru menskalakan
seluruh selangnya; dan jenazah yang dipindahkan dari tempat lain menyetel ulang
$T_a$ di tengah peluruhannya, sehingga memalsukan kematian yang lebih awal atau
lebih lambat. \omterm{def:g10:algebra:equation}{Persamaannya} jujur; anggapannyalah yang membawa risiko.

\textbf{17.} $z = \frac1y$ menuruti $z' = 1 - z$, jadi
$z = 1 + 9\eu^{-t}$ (dari $z(0) = 10$), sehingga
$y = \frac{1}{1 + 9\eu^{-t}}$. \omterm{def:g12:deriv:inflection}{Titik beloknya} di $y = \frac12$:
$9\eu^{-t} = 1$, jadi $t = \ln 9 \approx 2.2$ --- tanggal berbaliknya
wabah itu, langsung dari rumusnya.

\textbf{18.} Kecepatan terminalnya: $v' = 0$ ketika $v = 20$ m/s.
Penyelesaiannya: $v(t) = 20\left(1 - \eu^{-t/2}\right)$. Lalu untuk
$0.95$: $\eu^{-t/2} = 0.05$, jadi $t = 2\ln 20 \approx 6$ sekon ---
tetes hujan mencapai laju jelajahnya dalam hitungan sekon, dan itulah
sebabnya hujan tidak membunuh.

\textbf{19.} Keadaan mapannya: $c = 8$. Setengahnya:
$4 = 8\left(1 - \eu^{-t/2}\right)$, jadi $\eu^{-t/2} = \frac12$ dan
$t = 2\ln 2 \approx 1.4$ jam. Infusnya adalah kembaran sinambung
dari pinjaman itu: pemasukan tetap, pengeluaran sebanding ---
\cref{exo:g12:diffeq:9} membuat kamusnya menjadi eksak.

\textbf{20.} Peluruhan ($a < 0$, $b = 0$), pendinginan ($y = T - T_a$),
jatuh ($a < 0$, $b > 0$: mendekati kesetimbangan dari
bawah), penakaran obat (serupa, di dalam pembuluh): satu \omterm{def:g10:algebra:equation}{persamaan}, empat dunia.
Gelungnya: nyatakan hukum perubahannya; tulis \omterm{def:g10:algebra:equation}{persamaannya}; selesaikan
(\cref{met:g12:diffeq:solve}); tera tetapannya pada
pengukuran; ramalkan; lalu audit anggapannya (pertanyaan
16). Osilasi menuntut $y'' = -\omega^2 y$ --- yang dijumpai, diselesaikan, dan
diayunkan di \cref{pb:g12:trigo:1}.

\end{solution}
